\documentclass[11pt,a4paper]{article}
\usepackage[margin=1in]{geometry}
\usepackage{amsmath,amssymb,amsthm}
\usepackage{algorithm}
\usepackage{algpseudocode}
\usepackage{booktabs}
\usepackage{hyperref}

\theoremstyle{definition}
\newtheorem{definition}{Definition}[section]
\newtheorem{theorem}{Theorem}[section]
\newtheorem{lemma}[theorem]{Lemma}
\newtheorem{remark}{Remark}[section]

\title{\textbf{Rigorous Mathematical Specification, Channel Partitioning, and Spatial Statistical Dynamics of Group and Instance Normalization}}
\author{Team Scaibu \\ \small Email: \texttt{jaydeep.9664920749@gmail.com} \\ \small Architecture and Core Engine Team}
\date{October 2026}

\begin{document}
\maketitle

\begin{abstract}
This document presents the complete mathematical formulation, group partitioning mechanics, and computational characteristics of Group Normalization (GroupNorm) and Instance Normalization (InstanceNorm). Operating independently of the batch dimension, GroupNorm partitions the channel dimension into $G$ disjoint groups and normalizes activations over both the grouped channels and spatial locations for each individual sample. InstanceNorm represents the extreme case $G = C$, isolating channel-specific contrast. We detail the formal transformation, prove batch-size invariance, derive Big-O complexity bounds, trace a numerical example matching the reference implementation, and document numerical stability properties.
\end{abstract}

\section{Notation and Mathematical Preliminaries}

Let $X \in \mathbb{R}^{N \times C \times S}$ denote a 3D activation tensor representing $N$ samples, $C$ channels, and $S = H \times W$ spatial locations per channel.

\begin{itemize}
    \item $N \in \mathbb{N}_{\ge 1}$: Batch size.
    \item $C \in \mathbb{N}_{\ge 1}$: Number of feature channels.
    \item $S \in \mathbb{N}_{\ge 1}$: Number of spatial elements ($S = H \times W$).
    \item $G \in \mathbb{N}_{\ge 1}$: Number of groups such that $C \pmod G = 0$.
    \item $C_G = C/G$: Number of channels allocated to each group.
    \item $\mathcal{S}_g = \{c \in \mathbb{N} \mid g \cdot C_G \le c < (g+1) C_G\}$: Channel index set for group $g \in \{0, \dots, G-1\}$.
    \item $\mu_{n, g} \in \mathbb{R}$: Empirical mean of group $g$ for sample $n$.
    \item $\sigma_{n, g}^2 \in \mathbb{R}_{\ge 0}$: Empirical variance of group $g$ for sample $n$.
    \item $\gamma \in \mathbb{R}^C, \beta \in \mathbb{R}^C$: Per-channel learned scale and shift parameters.
    \item $\epsilon > 0$: Numerical regularization epsilon ($\epsilon = 10^{-5}$).
\end{itemize}

\begin{table}[h]
\centering
\caption{Summary of Mathematical Symbols for Group / Instance Normalization}
\begin{tabular}{lll}
\toprule
\textbf{Symbol} & \textbf{Type / Domain} & \textbf{Description} \\
\midrule
$X$ & $\mathbb{R}^{N \times C \times S}$ & 3D input activation tensor \\
$G$ & $\{1, \dots, C\}$ & Number of channel groups \\
$C_G$ & $\mathbb{N}_{\ge 1}$ & Channels per group ($C / G$) \\
$\mu_{n, g}$ & $\mathbb{R}$ & Mean for sample $n$ in group $g$ \\
$\sigma_{n, g}^2$ & $\mathbb{R}_{\ge 0}$ & Variance for sample $n$ in group $g$ \\
$\hat{X}$ & $\mathbb{R}^{N \times C \times S}$ & Standardized activation tensor \\
$Y$ & $\mathbb{R}^{N \times C \times S}$ & Normalized output tensor \\
$\gamma$ & $\mathbb{R}^C$ & Learned channel-wise scale vector \\
$\beta$ & $\mathbb{R}^C$ & Learned channel-wise shift vector \\
$\epsilon$ & $\mathbb{R}_{> 0}$ & Numerical stability constant \\
\bottomrule
\end{tabular}
\end{table}

\section{Problem Definition and Core Concepts}

\subsection{Intuitive Meaning}
Batch Normalization requires large batch sizes ($B \ge 32$) to compute stable population statistics, making it ineffective for high-resolution vision tasks (detection, segmentation, video generation) where hardware memory restricts batch sizes to $B \in \{1, 2\}$. Group Normalization addresses this limitation by computing statistics within subsets of channels per sample, making performance completely independent of batch size.

\subsection{Formal Mathematical Definition}
Given sample index $n \in \{1, \dots, N\}$, channel $c \in \{0, \dots, C-1\}$, and spatial location $s \in \{1, \dots, S\}$, let group index $g = \lfloor c / C_G \rfloor$. The group statistics are:
\begin{align}
\mu_{n, g} &= \frac{1}{C_G \cdot S} \sum_{k \in \mathcal{S}_g} \sum_{s=1}^S X_{n, k, s} \\
\sigma_{n, g}^2 &= \frac{1}{C_G \cdot S} \sum_{k \in \mathcal{S}_g} \sum_{s=1}^S (X_{n, k, s} - \mu_{n, g})^2 \\
\hat{X}_{n, c, s} &= \frac{X_{n, c, s} - \mu_{n, g}}{\sqrt{\sigma_{n, g}^2 + \epsilon}} \\
Y_{n, c, s} &= \gamma_c \hat{X}_{n, c, s} + \beta_c
\end{align}

\begin{remark}[Normalization Family Unification]
\begin{itemize}
    \item When $G = 1$: GroupNorm reduces to LayerNorm applied across all channels and spatial locations ($C_G = C$).
    \item When $G = C$: GroupNorm reduces to InstanceNorm ($C_G = 1$), normalizing each channel slice independently.
\end{itemize}
\end{remark}

\section{Algorithm Specification}

\begin{algorithm}[H]
\caption{Group Normalization Forward Pass}
\label{alg:group_norm}
\begin{algorithmic}[1]
\Require Activation tensor $X \in \mathbb{R}^{N \times C \times S}$, group count $G \in \mathbb{N}_{\ge 1}$, scale $\gamma \in \mathbb{R}^C$, shift $\beta \in \mathbb{R}^C$, regularizer $\epsilon > 0$.
\Ensure Normalized tensor $Y \in \mathbb{R}^{N \times C \times S}$, group means $\mu \in \mathbb{R}^{N \times G}$, group variances $\sigma^2 \in \mathbb{R}^{N \times G}$.
\State $N \gets \text{dim}_0(X)$, $C \gets \text{dim}_1(X)$, $S \gets \text{dim}_2(X)$
\State $C_G \gets C / G$, $M \gets C_G \cdot S$
\State Allocate $Y \in \mathbb{R}^{N \times C \times S}$, $\mu \in \mathbb{R}^{N \times G}$, $\sigma^2 \in \mathbb{R}^{N \times G}$
\For{$n \gets 1 \textbf{ to } N$}
    \For{$g \gets 0 \textbf{ to } G - 1$}
        \State $T_{\text{sum}} \gets 0$, $T_{\text{sq}} \gets 0$
        \For{$c \gets g \cdot C_G \textbf{ to } (g+1) C_G - 1$}
            \For{$s \gets 1 \textbf{ to } S$}
                \State $T_{\text{sum}} \gets T_{\text{sum}} + X_{n, c, s}$
            \EndFor
        \EndFor
        \State $\mu_{n, g} \gets T_{\text{sum}} / M$
        \For{$c \gets g \cdot C_G \textbf{ to } (g+1) C_G - 1$}
            \For{$s \gets 1 \textbf{ to } S$}
                \State $T_{\text{sq}} \gets T_{\text{sq}} + (X_{n, c, s} - \mu_{n, g})^2$
            \EndFor
        \EndFor
        \State $\sigma_{n, g}^2 \gets T_{\text{sq}} / M$
        \State $s_{\text{std}} \gets \sqrt{\sigma_{n, g}^2 + \epsilon}$
        \For{$c \gets g \cdot C_G \textbf{ to } (g+1) C_G - 1$}
            \For{$s \gets 1 \textbf{ to } S$}
                \State $\hat{X}_{n, c, s} \gets \frac{X_{n, c, s} - \mu_{n, g}}{s_{\text{std}}}$
                \State $Y_{n, c, s} \gets \gamma_c \hat{X}_{n, c, s} + \beta_c$
            \EndFor
        \EndFor
    \EndFor
\EndFor
\State \Return $\{Y, \mu, \sigma^2\}$
\end{algorithmic}
\end{algorithm}

\section{Theoretical Guarantees and Batch Invariance}

\begin{theorem}[Strict Batch Size Independence]
Let $X \in \mathbb{R}^{N \times C \times S}$ and let $X' \in \mathbb{R}^{N' \times C \times S}$ contain any subset of samples from $X$. For any sample $i$ present in both tensors:
\begin{equation}
\text{GN}(X')_i = \text{GN}(X)_i
\end{equation}
\end{theorem}

\begin{proof}
For any sample index $i$, the computation of group mean $\mu_{i, g}$ and variance $\sigma_{i, g}^2$ depends exclusively on the slice $X_{i, :, :}$. No terms contain cross-sample summations $\sum_{n=1}^N X_{n, c, s}$. Therefore, changing the batch dimension $N$ or the batch composition has zero effect on the output activations of sample $i$.
\end{proof}

\subsection{Limitations}
\begin{itemize}
    \item \textbf{Group Divisibility Constraint}: Number of channels $C$ must be an exact integer multiple of group count $G$.
\end{itemize}

\section{Complexity Analysis}

\subsection{Time Complexity}
\begin{itemize}
    \item \textbf{Group Statistics}: Computing mean and variance over $N \times C \times S$ takes $2 N C S$ additions.
    \item \textbf{Affine Scaling}: Applying $\gamma_c \hat{X} + \beta_c$ requires $3 N C S$ FLOPs.
    \item \textbf{Total Time Complexity}: $\Theta(N \cdot C \cdot S)$ FLOPs.
\end{itemize}

\subsection{Space Complexity}
\begin{itemize}
    \item \textbf{Auxiliary Memory}: Stores group means and variances of shape $(N, G)$, requiring $\mathcal{O}(N \cdot G)$ auxiliary memory.
\end{itemize}

\section{Exactness and Error Analysis}

The operation is mathematically exact. Utilizing two-pass accumulation for group variance guarantees numerical stability without subtractive cancellation.

\section{Worked Numerical Example}

Consider $N = 1, C = 2, S = 2, G = 1$ (hence $C_G = 2$, total elements $M = 4$):
\begin{equation}
X_{0, 0, :} = [1.0, 3.0], \quad X_{0, 1, :} = [5.0, 7.0]
\end{equation}
Parameters: $\gamma = [1.0, 1.0]^T, \beta = [0.0, 0.0]^T, \epsilon = 10^{-5}$.

\begin{enumerate}
    \item \textbf{Group Mean}:
    \begin{equation}
    \mu_{0, 0} = \frac{1.0 + 3.0 + 5.0 + 7.0}{4} = \frac{16.0}{4} = 4.0
    \end{equation}
    \item \textbf{Group Variance}:
    \begin{align}
    \sigma_{0, 0}^2 &= \frac{(1-4)^2 + (3-4)^2 + (5-4)^2 + (7-4)^2}{4} = \frac{9 + 1 + 1 + 9}{4} = \frac{20}{4} = 5.0
    \end{align}
    Standard deviation: $\sqrt{5.0 + 10^{-5}} = \sqrt{5.00001} \approx 2.2360679$.
    \item \textbf{Normalized Outputs}:
    \begin{align}
    Y_{0, 0, 0} &= \frac{1.0 - 4.0}{2.2360679} = \frac{-3.0}{2.2360679} \approx -1.34164 \\
    Y_{0, 0, 1} &= \frac{3.0 - 4.0}{2.2360679} = \frac{-1.0}{2.2360679} \approx -0.44721 \\
    Y_{0, 1, 0} &= \frac{5.0 - 4.0}{2.2360679} = \frac{1.0}{2.2360679} \approx 0.44721 \\
    Y_{0, 1, 1} &= \frac{7.0 - 4.0}{2.2360679} = \frac{3.0}{2.2360679} \approx 1.34164
    \end{align}
\end{enumerate}

\section{Numerical Considerations and Edge Cases}

\begin{itemize}
    \item \textbf{Uniform Constant Activations}: If all elements in a group are constant, variance is $0$; $\epsilon = 10^{-5}$ prevents division by zero.
    \item \textbf{Spatial Dimension $S = 1$}: Occurs in dense layers; GroupNorm transitions smoothly to grouped MLP normalization.
\end{itemize}

\begin{thebibliography}{9}
\bibitem{wu2018group} Y.~Wu and K.~He, ``Group Normalization,'' in \emph{Proceedings of the European Conference on Computer Vision (ECCV)}, pp.~3--19, 2018.
\bibitem{ulyanov2016instance} D.~Ulyanov, A.~Vedaldi, and V.~Lempitsky, ``Instance Normalization: The Missing Ingredient for Fast Stylization,'' \emph{arXiv preprint arXiv:1607.08022}, 2016.
\end{thebibliography}

\end{document}
